\section{Conclusiones y Trabajo Futuro}
\label{sec:conclusiones}
Tras completar el diseño del modelo en GAMA Platform, realizar las simulaciones experimentales bajo condiciones controladas y analizar los resultados cuantitativos obtenidos, se formularon las siguientes conclusiones:

\begin{itemize}
    \item Se diseñó de forma exitosa una representación virtual y espacial altamente realista de la infraestructura de red LAN universitaria. La integración directa de archivos vectoriales geográficos (Shapefiles) elaborados previamente en QGIS dentro del entorno de simulación GAMA demostró ser una solución técnica robusta para mapear de manera físicamente precisa la distribución geométrica de los activos y el cableado estructurado en los siete laboratorios de computación.
    \item Se modelaron satisfactoriamente los diversos dispositivos de red utilizando el paradigma de sistemas multiagente. La definición de comportamientos individuales y atributos específicos para cada especie (tales como listas de puertos TCP abiertos, estado booleano de infección, niveles de parcheo y roles de comunicación) permitió simular un entorno de red heterogéneo y dotar a cada host de respuestas dinámicas coherentes con un escenario real.
    \item Se simuló con precisión el vector de ingreso de la amenaza externa desde la WAN de Internet hacia la LAN interna. Las ejecuciones experimentales evidenciaron que, si bien la barrera perimetral provista por el firewall disminuyó de forma significativa la tasa de intrusión en los primeros ciclos de simulación, resultó insuficiente por sí sola para contener el ataque ante exposiciones continuas y repetidas solicitudes de conexión, consolidándose el brote una vez que la primera estación de trabajo fue comprometida.
    \item Se implementaron e integraron de manera satisfactoria reglas de propagación de carácter probabilístico basadas en el comportamiento histórico del ransomware WannaCry. La correlación del puerto SMB TCP 445 con variables de susceptibilidad y atenuación por parcheo (\texttt{patch\_}\allowbreak\texttt{level}) posibilitó recrear la dispersión lateral de forma autónoma a través de los enlaces físicos de adyacencia de la red.
    \item Se evaluó cuantitativamente el impacto destructivo del ataque informático sobre la infraestructura educativa universitaria. El escenario experimental E1 (propagación libre sin respuesta) reveló una vulnerabilidad crítica del campus, puesto que el ransomware comprometió de forma exponencial el $100\%$ de las estaciones de trabajo susceptibles en apenas 38 ciclos de simulación, dejando inoperables la totalidad de los laboratorios.
    \item Se validó experimentalmente la alta efectividad del protocolo de contingencia activa automatizada. La activación automática de medidas de aislamiento lógico y blindaje en hosts sanos al superar el umbral crítico del $30\%$ de equipos infectados (aproximadamente en el ciclo 18) contuvo de forma inmediata la dispersión lateral, protegiendo con éxito al $74,2\%$ restante de la red académica y demostrando que la respuesta rápida a nivel de terminal es la defensa más eficiente ante incidentes de gran escala.
\end{itemize}

\subsection{Líneas de Trabajo Futuro}
Como continuación de esta investigación, se proponen las siguientes líneas de desarrollo:
\begin{enumerate}
    \item \textbf{Segmentación Lógica de Redes (VLANs):} Incorporar al modelo la simulación de redes locales virtuales (VLANs) para evaluar el impacto de políticas de segmentación interna a nivel de switches, limitando la comunicación inter-laboratorios sin requerir el aislamiento completo de los hosts.
    \item \textbf{Modelado del Factor Humano:} Integrar agentes que representen a los usuarios de los laboratorios (estudiantes e investigadores) con reglas de comportamiento probabilísticas asociadas a la ingeniería social (por ejemplo, apertura de correos electrónicos de phishing como vector de entrada secundario).
    \item \textbf{Malware Multi-vector y Adaptativo:} Desarrollar reglas para simular ransomwares de última generación con múltiples capacidades de movimiento lateral (como el uso de credenciales comprometidas mediante herramientas de volcado de memoria en conjunto con exploits de red).
\end{enumerate}
